\section{The GEMM kernel}
\label{sec:kernel}

\subsection{Algorithm}
The kernel is \code{gemm\_hls}~\cite{gemmhls}, a communication-avoiding design for HLS. \Cref{fig:tiling} shows its
tiling. \(C\) is computed one \emph{memory tile} of \(T_N\times T_M\) values at a time; the tile stays on chip in the
processing elements (PEs) while the kernel streams the whole \(K\) dimension: at each step \(k\), column \(k\) of a
\(T_N\times K\) strip of \(A\) and row \(k\) of a \(K\times T_M\) strip of \(B\) form an outer product that is added to
the tile. A chain of \(P_N\) PEs, each \(P_M\) wide, performs \(P_N P_M\) multiply-adds per cycle, so the peak is
\(2 P_N P_M f\) flop/s. The cost model that explains most of the measurements in \cref{sec:eval}:
\begin{itemize}
  \item The time is \(NKM/(P_N P_M)\) cycles plus a fixed fill and drain cost per memory tile; that cost is amortised
        over \(K\), so efficiency rises with \(K\).
  \item The kernel always computes whole memory tiles: a matrix with \(N < T_N\) or \(M < T_M\) wastes the rest.
  \item \(A\) is read \(M/T_M\) times and \(B\) is read \(N/T_N\) times. The tile's arithmetic intensity,
        \(T_N T_M / (2(T_N+T_M))\) flop per byte of \(A\) and \(B\), is 32 for \(128^2\) tiles and 171 for
        \(1024\times512\): far above what a \SI{7.2}{GB/s} channel needs, so bandwidth is not the limit
        (\cref{sec:eval:roofline}).
  \item The tile lives in BRAM (\(T_N T_M\) values of C, double-buffered A). BRAM decides the largest tile.
\end{itemize}

\begin{figure}[t]
  \centering
  \resizebox{\linewidth}{!}{% Memory tiling: one T_N x T_M tile of C is computed from a T_N x K strip of A and a K x T_M strip of B.
\begin{tikzpicture}[x=1cm,y=1cm, every node/.append style={font=\sffamily\footnotesize}]
  % geometry: N = 4.0, K = 3.2, M = 4.8 (drawing units)
  \def\N{4.0} \def\K{3.2} \def\M{4.8} \def\tn{1.0} \def\tm{1.6}
  % B (K x M) on top right
  \draw[black!60, fill=cOrange!8] (0.6,0.6) rectangle (0.6+\M,0.6+\K);
  \fill[cOrange!40] (0.6+\tm,0.6) rectangle (0.6+2*\tm,0.6+\K);
  \node at (0.6+\M/2,0.6+\K+0.3) {B \quad \(K\times M\)};
  \draw[cOrange!70!black, thick] (0.6+\tm,0.6+\K*0.62) -- (0.6+2*\tm,0.6+\K*0.62);
  \node[lbl, text=cOrange!70!black, anchor=west] at (0.6+2*\tm+0.05,0.6+\K*0.62) {row \(k\): \(T_M\) values};
  % A (N x K) bottom left
  \draw[black!60, fill=cBlue!8] (-0.6-\K,-0.6-\N) rectangle (-0.6,-0.6);
  \fill[cBlue!35] (-0.6-\K,-0.6-2*\tn) rectangle (-0.6,-0.6-\tn);
  \node[rotate=90] at (-0.6-\K-0.35,-0.6-\N/2) {A \quad \(N\times K\)};
  \draw[cBlue!70!black, thick] (-0.6-\K*0.38,-0.6-2*\tn) -- (-0.6-\K*0.38,-0.6-\tn);
  \node[lbl, text=cBlue!70!black, anchor=north] at (-0.6-\K*0.38,-0.6-2*\tn-0.02) {column \(k\): \(T_N\) values};
  % C (N x M) bottom right, tiled
  \draw[black!60, fill=cGreen!5] (0.6,-0.6-\N) rectangle (0.6+\M,-0.6);
  \foreach \i in {1,2,3} { \draw[black!25] (0.6,-0.6-\i*\tn) -- (0.6+\M,-0.6-\i*\tn); }
  \foreach \j in {1,2} { \draw[black!25] (0.6+\j*\tm,-0.6) -- (0.6+\j*\tm,-0.6-\N); }
  \fill[cGreen!45] (0.6+\tm,-0.6-2*\tn) rectangle (0.6+2*\tm,-0.6-\tn);
  \node at (0.6+1.5*\tm,-0.6-1.5*\tn) {\(T_N\times T_M\)};
  \node at (0.6+\M/2,-0.6-\N-0.75) {C \quad \(N\times M\), computed one memory tile at a time};
  \draw[decorate, decoration={brace, amplitude=3pt}, black!60] (0.6+\M+0.08,-0.6-\tn) -- node[right=3pt, lbl] {\(T_N\)} (0.6+\M+0.08,-0.6-2*\tn);
  \draw[decorate, decoration={brace, amplitude=3pt, mirror}, black!60] (0.6+\tm,-0.6-\N-0.08) -- node[below=3pt, lbl] {\(T_M\)} (0.6+2*\tm,-0.6-\N-0.08);
  % outer product arrows
  \draw[->, thick, cBlue!70!black] (-0.6,-0.6-1.5*\tn) -- (0.6+\tm-0.05,-0.6-1.5*\tn);
  \draw[->, thick, cOrange!70!black] (0.6+1.5*\tm,0.6) -- (0.6+1.5*\tm,-0.6-\tn+0.05);
  % explanation
  \node[blk, fill=white, align=left, text width=70mm, font=\sffamily\scriptsize, anchor=north west] at (6.3,3.75)
   {\textbf{For each memory tile} (\(N/T_N \times M/T_M\) of them):\\[2pt]
    for \(k = 0 \ldots K{-}1\): C tile \(\mathrel{+}=\) (column \(k\) of the A strip) \(\otimes\) (row \(k\) of the B strip),
    an outer product of \(T_N T_M\) multiply-adds in \(\dfrac{T_N T_M}{P_N P_M}\) cycles.\\[3pt]
    \textbf{Time:} \(\dfrac{N K M}{P_N P_M}\) cycles + a fixed fill/drain cost per tile (load the first A slice,
    drain \(T_N T_M\) values of C through the PE chain), amortised over \(K\).\\[3pt]
    \textbf{Traffic:} A is read \(M/T_M\) times, B \(N/T_N\) times; C is written once.\\[3pt]
    \textbf{Intensity} (A + B, float32): \(\dfrac{2 T_N T_M K}{4K(T_N+T_M)} = \dfrac{T_N T_M}{2(T_N+T_M)}\) flop/byte,
    e.g.\ 32 for \(128^2\) tiles and 171 for \(1024\times512\).\\[3pt]
    \textbf{On chip:} C tile in the PEs' BRAM (\(T_N T_M\) values); this is what limits the tile size (BRAM).};
\end{tikzpicture}
}
  \caption{Memory tiling of \code{gemm\_hls}. The highlighted tile of \(C\) is computed from one strip of \(A\) and
    one strip of \(B\), one outer product per step of \(k\).}
  \label{fig:tiling}
\end{figure}

\subsection{Dataflow}
\Cref{fig:kernel} shows the hardware as HLS builds it: concurrent functions connected by FIFOs. Readers fetch A and B
in bursts and convert 512-bit memory words into vectors for the PEs. \(A\) values and \(B\) vectors travel forward
along the PE chain; each PE keeps its rows of the C tile and a double-buffered slice of A, so that the next outer
product can be loaded while the current one is computed. When a tile is finished, its values drain backwards through
the chain to PE~0, are packed into 512-bit words and written.

\begin{figure}[t]
  \centering
  \resizebox{\linewidth}{!}{% The gemm_hls dataflow (transposed-A variant): memory readers, width converters, PE chain, C drain and writer.
\begin{tikzpicture}[x=1cm,y=1cm, every node/.append style={font=\sffamily\footnotesize}]
  \tikzset{fn/.style={acc, minimum width=24mm, minimum height=9mm},
           pe/.style={acc, fill=cOrange!32, minimum width=12mm, minimum height=16mm, align=center},
           mem/.style={ddr, minimum width=14mm, minimum height=9mm}}
  % memory side
  \node[mem] (ma) at (0, 2.2)  {\code{gmem0}\\[-1pt]\scriptsize A\(^\mathsf{T}\)};
  \node[mem] (mb) at (0, 0)    {\code{gmem1}\\[-1pt]\scriptsize B};
  \node[mem] (mc) at (0,-2.4)  {\code{gmem2}\\[-1pt]\scriptsize C};
  \node[fn] (ra) at (2.9, 2.2)  {\code{ReadATransposed}\\[-1pt]\scriptsize 16-beat bursts};
  \node[fn] (ca) at (6.0, 2.2)  {\code{\scriptsize ConvertWidthATransposed}\\[-1pt]\scriptsize 512\,b \(\rightarrow\) \(P_N\) values};
  \node[fn] (rb) at (2.9, 0)    {\code{ReadB}\\[-1pt]\scriptsize row bursts};
  \node[fn] (fb) at (6.0, 0)    {\code{FeedB}\\[-1pt]\scriptsize \(P_M\) values / cycle};
  \node[fn] (wc) at (2.9,-2.4)  {\code{WriteC}};
  \node[fn] (cc) at (6.0,-2.4)  {\code{ConvertWidthC}\\[-1pt]\scriptsize \(P_M\) values \(\rightarrow\) 512\,b};
  \draw[wbus] (ma) -- (ra);  \draw[wbus] (ra) -- (ca);
  \draw[wbus] (mb) -- (rb);  \draw[wbus] (rb) -- (fb);
  \draw[wbus] (cc) -- (wc);  \draw[wbus] (wc) -- (mc);
  % PE chain
  \foreach \i/\x/\t in {0/8.9/0, 1/10.5/1, 2/12.1/2, 3/14.4/{P_N{-}1}} {
    \node[pe] (pe\i) at (\x, 0) {PE\\[-1pt]\(\t\)\\[2pt]\scriptsize \(P_M\) MAC};
  }
  \node[font=\large] at (13.25,0) {\(\cdots\)};
  \draw[link, cBlue!80!black] (ca.east) -| node[lbl, pos=0.25, above, text=cBlue!80!black] {aPipes[0]} (pe0.north);
  \draw[link, cBlue!80!black] (fb) -- node[lbl, above, text=cBlue!80!black] {bPipes[0]} (pe0);
  \foreach \i/\j in {0/1, 1/2} {
    \draw[link, cBlue!80!black] ([yshift=4mm]pe\i.east) -- ([yshift=4mm]pe\j.west);
    \draw[link, cBlue!80!black] ([yshift=0mm]pe\i.east) -- ([yshift=0mm]pe\j.west);
    \draw[link, cGreen!60!black] ([yshift=-4mm]pe\j.west) -- ([yshift=-4mm]pe\i.east);
  }
  \draw[link, cBlue!80!black] ([yshift=4mm]pe2.east) -- ++(0.4,0);
  \draw[link, cBlue!80!black] ([yshift=0mm]pe2.east) -- ++(0.4,0);
  \draw[link, cGreen!60!black] ([xshift=-0.4cm,yshift=-4mm]pe3.west) -- ([yshift=-4mm]pe2.east);
  \draw[link, cBlue!80!black] ([xshift=-0.4cm,yshift=4mm]pe3.west) -- ([yshift=4mm]pe3.west);
  \draw[link, cBlue!80!black] ([xshift=-0.4cm]pe3.west) -- (pe3.west);
  \draw[link, cGreen!60!black] ([yshift=-4mm]pe3.west) -- ++(-0.4,0);
  \draw[link, cGreen!60!black] (pe0.south) |- node[lbl, pos=0.75, above, text=cGreen!60!black] {cPipes[0]} (cc.east);
  % legend for the pipes
  \node[note, anchor=west, align=left] (lg) at (9.4,-1.75)
     {\textcolor{cBlue!80!black}{\rule[2pt]{6mm}{1pt}} A and B values move forward along the chain\\
      \textcolor{cGreen!60!black}{\rule[2pt]{6mm}{1pt}} finished C tiles drain backwards to PE 0};
  \node[note, anchor=west, align=left, text width=58mm] (nt) at (9.4,-2.75)
     {Each PE keeps \(\frac{T_N}{P_N}\times T_M\) values of C (\code{cBuffer}, BRAM) and a double-buffered slice of A.
      All functions run concurrently (\code{\#pragma HLS DATAFLOW}), connected by FIFOs.};
  \begin{scope}[on background layer]
    \node[region=cOrange, fit=(ma)(mc)(pe3)(lg)(nt), inner sep=6pt] (ip) {};
  \end{scope}
  \node[regionl=cOrange, anchor=south west] at (ip.north west)
     {\code{MatrixMultiplicationKernel} (\file{kernels/gemm/src/Top.cpp}), accel clock, \code{ap\_ctrl\_hs}};
\end{tikzpicture}
}
  \caption{Kernel dataflow (transposed-A variant). All functions run at the same time; the C values of a finished
    tile drain backwards through the PE chain.}
  \label{fig:kernel}
\end{figure}

\paragraph{Transposed A.} Without \code{transposed\_a}, the kernel reads \(A\) row-major and transposes it on chip,
which means reading single 64-byte words with a stride: one AXI burst per word. With \code{transposed\_a} the host
sends \(A^\mathsf{T}\) and the kernel reads \(A\) in 16-beat bursts. For PE 32\x8 this was the bottleneck
(\cref{sec:eval:layout}).

\subsection{Changes to gemm\_hls}
The upstream repository stays unmodified as a submodule; \file{kernels/gemm/} is a copy, and every change is a
separate commit listed in \file{kernels/gemm/ORIGIN.md}:
\begin{itemize}
  \item \textbf{Interfaces for a shell-less SoC.} \code{ap\_ctrl\_hs} block control, \code{m\_axi offset=direct}
        (each matrix address becomes an input port) and \code{ap\_none} sizes, with no AXI-Lite: the SoC drives
        them from CSRs. Burst length, outstanding requests and latency of the masters come from the config.
  \item \textbf{Vitis HLS 2026.1.} Arrays of \code{hlslib::Stream} objects fail in synthesis (HLS 214-177, unresolved
        Phi node), so synthesis uses \code{hls::stream} while simulation keeps hlslib~\cite{hlslib}.
        \code{ConvertWidthC} indexed a 512-bit word dynamically in an II=1 loop (estimated 10.6\,ns); it now
        unpacks a word with constant indices.
  \item \textbf{Knobs.} FIFO depths and all tiling parameters come from \code{Config.h}, which the flow renders
        from the JSON config.
\end{itemize}

\subsection{Integration into the SoC}
\code{HlsGemmKernel} (\file{hw/gemm/hls\_kernel.py}, \cref{fig:integration}) instantiates the exported IP. Its port
list is generated from the HLS report (\code{csynth.xml}), so no port name or width is written by hand.
\begin{itemize}
  \item \textbf{Control across clock domains.} The \code{start} CSR becomes a pulse in the accel domain; a small
        state machine holds \code{ap\_start} until \code{ap\_ready} and counts cycles until \code{ap\_done}. The
        arguments are written before the start and do not change while the kernel runs, so a MultiReg per bit is
        enough; status and counters cross back the same way and are read after \code{done}.
  \item \textbf{AXI to LiteDRAM.} A and B use LiteDRAM's AXI-to-native read converter. For C the project has its own
        \code{AxiWriteToNative} (\file{hw/gemm/axi\_native.py}): with LiteDRAM's write converter, Vivado removed the read
        side of a response FIFO whose content was constant but kept the cell, leaving driverless pins
        (\code{Route 35-13}). Each converter is built for one direction only, for the same reason.
  \item \textbf{Performance counters.} \code{AxiPerfMonitor} (\file{hw/gemm/axi\_perf.py}) watches each master while
        the kernel runs: handshakes, beats, address stalls, cycles with the read channel ready but no data
        (\code{rready \& !rvalid}, called \emph{read waits} below), back-pressure, write stalls, and the read latency
        from each AR handshake to its first R beat (valid because the HLS masters use a single ID, so data returns in
        order).
  \item \textbf{Vivado import.} A Tcl hook before synthesis (\file{tcl/vivado/import\_hls\_ip.tcl}) adds the HLS IP
        repository, creates the IP and synthesises it, including its floating-point cores.
\end{itemize}

\begin{figure}[t]
  \centering
  \resizebox{\linewidth}{!}{% How the HLS IP is wired into the SoC: CSR control across clock domains, AXI to LiteDRAM native, perf taps.
\begin{tikzpicture}[x=1cm,y=1cm, every node/.append style={font=\sffamily\footnotesize}]
  \tikzset{w/.style={minimum width=#1}, cs/.style={sys, minimum width=24mm, minimum height=6mm, font=\sffamily\scriptsize}}
  % CSRs (sys)
  \node[cs] (c1) at (0, 1.5)  {\code{start} (write pulse)};
  \node[cs] (c2) at (0, 0.75) {\code{a, b, c, size\_n/k/m}};
  \node[cs] (c3) at (0, 0)    {\code{soft\_reset}};
  \node[cs] (c4) at (0,-0.75) {\code{status} (done/idle/busy)};
  \node[cs] (c5) at (0,-1.5)  {\code{cycles} (64 b)};
  \node[cs] (c6) at (0,-2.25) {\code{perf\_\{a,b,c\}\_*}};
  \node[cs] (c7) at (0,-3.0)  {\code{build\_id}};
  \begin{scope}[on background layer]\node[region=cBlue, fit=(c1)(c7), inner sep=5pt] (csr) {};\end{scope}
  \node[regionl=cBlue, anchor=south] at (csr.north) {CSRs, sys};
  % CDC
  \node[blk, fill=cRed!8, draw=cRed!70!black, w=21mm, font=\sffamily\scriptsize] (ps) at (3.1, 1.5)  {PulseSynchronizer};
  \node[blk, fill=cRed!8, draw=cRed!70!black, w=21mm, font=\sffamily\scriptsize] (mr) at (3.1, 0.38) {MultiReg};
  \node[blk, fill=cRed!8, draw=cRed!70!black, w=21mm, font=\sffamily\scriptsize] (mb) at (3.1,-1.5)  {MultiReg};
  \draw[link] (c1) -- (ps);
  \draw[link] (c2.east) -- ++(0.3,0) |- ([yshift=1mm]mr.west);
  \draw[link] (c3.east) -- ++(0.3,0) |- ([yshift=-1mm]mr.west);
  \draw[link] (mb.west) -- ++(-0.3,0) |- (c4.east);
  \draw[link] (mb) -- (c5);
  \draw[link] (mb.west) -- ++(-0.3,0) |- (c6.east);
  % control FSM (accel)
  \node[acc, w=27mm, align=left, font=\sffamily\scriptsize, text width=25mm] (fsm) at (6.2,0.6)
       {\textbf{control} (accel)\\ hold \code{ap\_start} until \code{ap\_ready}; count cycles while busy; done on \code{ap\_done}};
  \draw[link] (ps) -| (fsm.north);
  \draw[link] (mr) -- (mr -| fsm.west);
  \draw[link] (fsm.south) |- (mb.east);
  % HLS IP
  \node[acc, fill=cOrange!32, w=30mm, minimum height=40mm, align=center] (ip) at (10.0,0.0)
       {\textbf{HLS IP}\\ \code{gemm\_kernel\_ip}\\[4pt]\scriptsize Instance generated from\\ \scriptsize the \code{csynth.xml} port list\\[4pt]
        \scriptsize \code{ap\_ctrl\_hs}, \code{ap\_none} args,\\ \scriptsize \code{m\_axi offset=direct}};
  \draw[link] (fsm.east |- 0,1.2) -- node[lbl, above] {start / args} (ip.west |- 0,1.2);
  \draw[link] (ip.west |- 0,0.0) -- node[lbl, above] {done, idle} (fsm.east |- 0,0.0);
  % AXI paths
  \foreach \m/\y/\conv in {A/1.4/\code{LiteDRAMAXI2NativeR}, B/0.0/\code{LiteDRAMAXI2NativeR}, C/-1.4/\code{AxiWriteToNative}} {
    \node[acc, w=8mm, minimum height=6mm, font=\sffamily\scriptsize] (tap\m) at (12.6,\y) {perf};
    \node[acc, w=30mm, font=\sffamily\scriptsize] (cv\m) at (15.0,\y) {\conv};
    \node[sys, w=16mm, font=\sffamily\scriptsize] (np\m) at (17.75,\y) {native port};
    \draw[wbus] (ip.east |- tap\m) -- node[lbl, above, pos=0.45] {AXI4 \m} (tap\m);
    \draw[wbus] (tap\m) -- (cv\m);
    \draw[wbus] (cv\m) -- (np\m);
  }
  \node[note, anchor=north, text width=60mm] at (15.4,-2.0)
     {Counters per bundle: AR/AW handshakes, R/W beats, address stalls, cycles waiting for read data
      (\code{rready \& !rvalid}), back-pressure, write stalls, read latency min/max/sum (AR \(\rightarrow\) first R, single ID).};
  \draw[ctl] (tapC.south) |- ++(-1.5,-1.85) -| node[lblw, pos=0.25, below] {counters (64 b each)} (mb.south);
  \node[note, anchor=north, text width=40mm] at (6.2,-2.1)
     {Arguments are quasi-static: written before \code{start}, so a MultiReg per bit is safe. Counters are read after \code{done}, when they are stable.};
\end{tikzpicture}
}
  \caption{The kernel inside the SoC: CSR control with clock-domain crossings, the AXI masters converted to LiteDRAM
    native ports, and a performance monitor on each master.}
  \label{fig:integration}
\end{figure}
